% !TeX program = xelatex
\documentclass[aspectratio=169,10pt]{beamer}

\usepackage{polyglossia}
\setdefaultlanguage{russian}

\usetheme{Madrid}
\usecolortheme{default}
\setbeamertemplate{navigation symbols}{}
\setbeamertemplate{footline}[frame number]

\usepackage{graphicx}
\usepackage{tikz}
\usetikzlibrary{arrows.meta,positioning,fit,shapes.geometric}
\usepackage{booktabs}
\usepackage{tabularx}
\usepackage{array}
\usepackage{multirow}
\usepackage{ragged2e}
\usepackage{makecell}
\usepackage{xcolor}
\usepackage{hyperref}
\usepackage{listings}
\usepackage{caption}

\hypersetup{
  colorlinks=true,
  linkcolor=blue,
  urlcolor=blue
}

\definecolor{codebg}{RGB}{248,248,248}
\definecolor{codeframe}{RGB}{210,210,210}
\definecolor{darkgreen}{RGB}{0,110,0}
\definecolor{darkblue}{RGB}{0,60,140}

\lstdefinelanguage{PlantUML}{
  morekeywords={
    @startuml,@enduml,title,actor,participant,database,rectangle,component,node,
    artifact,cloud,folder,package,frame,queue,collections,note,left,right,of,as,
    start,stop,if,then,else,endif,repeat,while,endwhile,skinparam
  },
  sensitive=false,
  morecomment=[l]',
  morestring=[b]"
}

\lstset{
  language=PlantUML,
  basicstyle=\ttfamily\scriptsize,
  keywordstyle=\color{darkblue}\bfseries,
  commentstyle=\color{darkgreen},
  stringstyle=\color{red!60!black},
  breaklines=true,
  columns=fullflexible,
  frame=single,
  rulecolor=\color{codeframe},
  backgroundcolor=\color{codebg},
  showstringspaces=false,
  tabsize=2
}
\input{lk-theme.tex}

% Кириллица в Beamer/TikZ: явные шрифты Unicode (иначе LM Sans без кириллицы в узлах)
\IfFontExistsTF{DejaVu Sans}{%
  \setsansfont{DejaVu Sans}%
  \setmainfont{DejaVu Serif}%
  \setmonofont{DejaVu Sans Mono}%
  \newfontfamily\cyrillicfont{DejaVu Serif}%
  \newfontfamily\cyrillicfontsf{DejaVu Sans}%
  \newfontfamily\cyrillicfonttt{DejaVu Sans Mono}%
}{%
  \setsansfont{Latin Modern Sans}%
  \setmainfont{Latin Modern Roman}%
  \setmonofont{Latin Modern Mono}%
}
\usefonttheme{professionalfonts}

% Пути к рисункам (сборка из каталога docs/slides/sbom или с -output-directory)
\graphicspath{{images/}{./images/}{docs/slides/sbom/images/}}

\newcolumntype{Y}{>{\RaggedRight\arraybackslash}X}

\title[SBOM]{Методически корректное формирование SBOM}
\subtitle{Подход, процессы, контроль качества и интеграция в жизненный цикл ПО}
\author[Методика SBOM]{Презентация для методического и практического использования}
\date{\today\quad(ГОСТ~Р~72118-2025, TARA, ДВБ)}

\begin{document}

\begin{frame}
  \titlepage
\end{frame}

\begin{frame}{План выступления}
\tableofcontents
\end{frame}

\section{Зачем нужен SBOM}

\begin{frame}{Почему тема SBOM стала обязательной}
\begin{itemize}
  \item Рост числа атак на цепочку поставки ПО.
  \item Повышение требований заказчиков, регуляторов и аудита.
  \item Необходимость быстро отвечать на вопросы:
  \begin{itemize}
    \item что входит в поставку;
    \item где используется уязвимый компонент;
    \item какие лицензии и обязательства применимы.
  \end{itemize}
  \item Переход от инвентаризации «по памяти команды» к формальному артефакту.
\end{itemize}

\begin{block}{Ключевая мысль}
SBOM — это не выгрузка из сканера, а управляемый артефакт жизненного цикла продукта.
\end{block}
\end{frame}

\begin{frame}{Что такое SBOM}
\begin{block}{Определение}
SBOM (Software Bill of Materials) — это формализованный машиночитаемый перечень компонентов, входящих в программный продукт, сборку, образ, пакет или поставку.
\end{block}

\begin{itemize}
  \item отражает состав продукта;
  \item фиксирует версии и зависимости;
  \item обеспечивает прослеживаемость происхождения компонентов;
  \item служит входом для процессов безопасности, лицензирования и аудита.
\end{itemize}
\end{frame}

\begin{frame}{Каким должен быть корректный SBOM}
\begin{table}
\footnotesize
\begin{tabularx}{\textwidth}{>{\bfseries}p{2.7cm}Y}
\toprule
Свойство & Пояснение \\
\midrule
Полный & Охватывает все значимые компоненты поставки, а не только прямые зависимости. \\
Прослеживаемый & Позволяет понять источник данных: манифест, lock-файл, бинарный анализ, контейнер, реестр артефактов. \\
Актуальный & Привязан к конкретной сборке, релизу, образу или поставке. \\
Машиночитаемый & Представлен в стандартном формате и пригоден для автоматической обработки. \\
Повторяемый & Для одинакового процесса и одинакового входа дает воспроизводимый результат. \\
Проверяемый & Допускает валидацию схемы, полноты, идентификации и качества. \\
\bottomrule
\end{tabularx}
\end{table}
\end{frame}
\begin{frame}{Основные цели формирования SBOM}
\begin{columns}[T]
\begin{column}{0.5\textwidth}
\begin{block}{Безопасность}
\begin{itemize}
  \item поиск уязвимых компонентов;
  \item реакция на CVE;
  \item мониторинг цепочки поставки;
  \item поддержка VEX.
\end{itemize}
\end{block}
\end{column}
\begin{column}{0.5\textwidth}
\begin{block}{Управление и комплаенс}
\begin{itemize}
  \item лицензионный контроль;
  \item выполнение контрактных требований;
  \item аудит релизов;
  \item унификация релизного процесса.
\end{itemize}
\end{block}
\end{column}
\end{columns}

\begin{alertblock}{Практический вывод}
Методика построения SBOM зависит от цели. Если цель не определена, артефакт быстро теряет практическую ценность.
\end{alertblock}
\end{frame}

\section{Объект учета и границы}

\begin{frame}{Для чего именно формируется SBOM}
\begin{itemize}
  \item для исходного проекта;
  \item для конкретной сборки;
  \item для контейнерного образа;
  \item для инсталляционного пакета;
  \item для встроенного ПО;
  \item для комплексного решения из нескольких сервисов.
\end{itemize}

\begin{block}{Что должно быть зафиксировано методически}
Единица учета, границы состава, уровень детализации, момент формирования и ответственный за выпуск.
\end{block}
\end{frame}

\begin{frame}{Типовые вопросы, которые надо решить заранее}
\begin{itemize}
  \item Включать ли транзитивные зависимости?
  \item Нужно ли учитывать компоненты базового контейнерного образа?
  \item Нужно ли включать системные пакеты ОС?
  \item Включать ли плагины, скрипты, модели, данные конфигурации?
  \item Нужно ли отражать инструменты сборки?
  \item Как учитывать внутренние проприетарные модули и форки?
\end{itemize}

\begin{alertblock}{Риск}
Если границы состава не формализованы, разные команды начнут выпускать несопоставимые SBOM.
\end{alertblock}
\end{frame}

\begin{frame}{Типы SBOM по моменту формирования}
\begin{table}
\scriptsize
\setlength{\tabcolsep}{3pt}
\begin{tabularx}{\textwidth}{>{\bfseries}p{2.2cm}p{2.85cm}Y}
\toprule
Тип & Источник & Для чего полезен \\
\midrule
Design-time &
\makecell[tl]{Манифесты, lock-файлы,\\описание проекта} &
Анализ заявленных зависимостей на стадии разработки. \\
Build-time &
\makecell[tl]{Сборка, пакеты, артефакты,\\контейнерные слои} &
Фиксация того, что реально вошло в сборку и релиз. \\
Runtime / Deployed &
\makecell[tl]{Среда исполнения, контейнер,\\ОС-пакеты} &
Проверка фактически развёрнутого состояния и дрейфа от релиза. \\
\bottomrule
\end{tabularx}
\end{table}

\begin{block}{Рекомендуемый подход}
Минимум — иметь build-time SBOM, а для зрелой практики дополнять его design-time и runtime.
\end{block}
\end{frame}

\section{Стандарты и источники данных}

\begin{frame}{Сравнение форматов: SPDX и CycloneDX}
\begin{table}
\tiny
\setlength{\tabcolsep}{3pt}
\renewcommand{\arraystretch}{1.12}
\begin{tabularx}{\textwidth}{>{\bfseries}p{1.55cm}Y Y Y}
\toprule
Формат & Сильные стороны & Когда выбирать & Ограничения \\
\midrule
SPDX &
\makecell[tl]{Зрелый стандарт; сильная\\лицензионная модель} &
\makecell[tl]{Юридический и лицензионный\\учёт; обмен с заказчиком} &
\makecell[tl]{В части DevSecOps-сценариев\\менее удобен, чем CycloneDX} \\
CycloneDX &
\makecell[tl]{Удобен для безопасности,\\зависимостей, контейнеров, VEX} &
\makecell[tl]{DevSecOps; уязвимости; CI/CD} &
\makecell[tl]{Нужна дисциплина по ID\\и метаданным} \\
\bottomrule
\end{tabularx}
\end{table}
\end{frame}

\begin{frame}{Сравнение форматов: SWID}
\begin{table}
\footnotesize
\begin{tabularx}{\textwidth}{>{\bfseries}p{1.8cm}Y Y Y}
\toprule
Формат & Сильные стороны & Когда выбирать & Ограничения \\
\midrule
SWID &
\makecell[tl]{Учёт установленного ПО\\в корпоративных сценариях} &
\makecell[tl]{Специализированные\\корпоративные контуры} &
\makecell[tl]{Реже — основной формат SBOM} \\
\bottomrule
\end{tabularx}
\end{table}
\vspace{0.4em}
\footnotesize
Выбор формата закрепляют в политике; при необходимости допускают два формата (например SPDX + CycloneDX) для разных контуров.
\end{frame}

\begin{frame}{Из каких источников должен строиться SBOM}
\begin{table}
\scriptsize
\begin{tabularx}{\textwidth}{>{\bfseries}p{3.0cm}Y}
\toprule
Источник & Примеры \\
\midrule
Манифесты зависимостей & pom.xml, package-lock.json, requirements.txt, poetry.lock, go.mod, Cargo.lock. \\
Сборочная система & Maven, Gradle, npm, pip, Poetry, Go, Cargo, Bazel. \\
Бинарный анализ & Сканирование библиотек, архивов, jar, dll, so, исполняемых файлов. \\
Контейнерный анализ & Образы, слои, базовые образы, ОС-пакеты внутри контейнера. \\
Реестр артефактов & Репозитории пакетов, контейнерные регистри, хранилища релизов. \\
Ручное подтверждение & Внутренние компоненты, форки, подрядные модули, патченные сборки. \\
\bottomrule
\end{tabularx}
\end{table}

\begin{alertblock}{Правильный принцип}
Нужно учитывать не только то, что заявлено в проекте, но и то, что реально попало в поставку.
\end{alertblock}
\end{frame}

\section{ГОСТ Р 72118-2025, TARA и доверенная вычислительная база}

\begin{frame}{ГОСТ Р 72118-2025: официальное издание (по PDF)}
\footnotesize
\textbf{Полное наименование:} Защита информации. Системы с конструктивной информационной безопасностью. \textbf{Методология разработки} (ГОСТ Р 72118---2025).

\vspace{0.25em}
\begin{itemize}
  \setlength{\itemsep}{0.15em}
  \item \textbf{Разработан:} ФСТЭК России; ФАУ «ГНИИИ ПТЗИ ФСТЭК России»; АО «Лаборатория Касперского»; ИСП РАН.
  \item \textbf{Внесён:} ТК~362 «Защита информации».
  \item \textbf{Утверждён и введён в действие:} Приказ Росстандарта от 6~июня~2025~г. №~539-ст (издание официальное, 2025).
\end{itemize}

\end{frame}

\begin{frame}{Проект стандарта (DOCX) и связь с SBOM}

\vspace{0.35em}
\textbf{Связь с SBOM} (инженерная интерпретация для поставки ПО):
\begin{itemize}
  \setlength{\itemsep}{0.15em}
  \item Стандарт задаёт \textbf{методологию СКИБ} и опирается на требования по ЖЦ (ГОСТ~Р~57193-2016) и на \textbf{безопасную разработку ПО} (ГОСТ~Р~56939) — SBOM фиксирует \textbf{фактический состав} компонентов, подпадающих под эти процессы.
  \item Для обоснования \textbf{целей безопасности} и снижения последствий угроз (в т.\,ч.\ эксплуатация уязвимостей) важны уменьшение поверхности атаки и контроль связей между элементами — \textbf{прозрачный SBOM} это поддерживает.
  \item Форматы CycloneDX/SPDX остаются практической основой машиночитаемого SBOM; ГОСТ определяет \textbf{организационно-методический контекст} применения.
\end{itemize}
\end{frame}

\begin{frame}{Термины ГОСТ Р 72118-2025 (§\,3, часть 1)}
\scriptsize
\begin{description}
  \setlength{\itemsep}{0.1em}
  \item[СКИБ] система с конструктивной ИБ (достигнутой конструктивными подходами).
  \item[ЦБ] намерение обеспечить заданные характеристики безопасности; проверяемость по критериям.
  \item[ПБ] ограничения и допущения среды, дополняющие ЦБ.
\end{description}
\end{frame}

\begin{frame}{Термины ГОСТ Р 72118-2025 (§\,3, часть 2)}
\scriptsize
\begin{description}
  \setlength{\itemsep}{0.1em}
  \item[Доверенная система] обоснованное соответствие ЦБ при выполнении ПБ (доверенным может быть элемент).
  \item[Конструктивный подход] свойства безопасности закладываются при создании, с замысла (архитектура, шаблоны).
\end{description}
\vspace{0.35em}
\footnotesize \textbf{ДВБ} на слайдах — инженерный аналог минимального доверенного набора компонентов (см. доверенную систему в стандарте).
\end{frame}

\begin{frame}{Методология СКИБ и ЖЦ (§\,5): процессы}
\footnotesize
Методология СКИБ встраивается в модель ЖЦ (ГОСТ~Р~72118-2025) и охватывает процессы:
\begin{itemize}
  \setlength{\itemsep}{0.08em}
  \item планирование; анализ бизнеса\/назначения; требования заинтересованных сторон;
  \item \textbf{системные требования}; \textbf{архитектура}; проект; системный анализ;
  \item \textbf{реализация}; комплексирование; гарантия качества; \textbf{верификация}; \textbf{валидация}.
\end{itemize}
\end{frame}

\begin{frame}{Методология СКИБ: где SBOM}
\footnotesize
\begin{itemize}
  \setlength{\itemsep}{0.15em}
  \item При \textbf{реализации и комплексировании} SBOM фиксирует фактический состав поставляемого ПО.
  \item На \textbf{верификации и валидации} SBOM сопоставляют с требованиями, ЦБ доменов и результатами анализа уязвимостей.
  \item Параллельно с \textbf{системной инженерией} цели SBOM (полнота, границы, форматы) задаются на уровне политики и трассируются к ЦБ.
\end{itemize}
\end{frame}

\begin{frame}{Трассировка: моделирование угроз $\rightarrow$ цели $\rightarrow$ ДВБ $\rightarrow$ SBOM}
\footnotesize
\textbf{TARA} задаёт сценарии и пути атаки к активам. В терминологии ГОСТ~Р~72118-2025 \textbf{цели безопасности (ЦБ)} формализуют намерение по защите; \textbf{предположения безопасности (ПБ)} задают контекст среды. \textbf{ДВБ} здесь — инженерное обозначение минимального доверенного набора ПО/аппаратуры (см. также «доверенная система» в стандарте). SBOM описывает \textbf{фактический состав} компонентов внутри и на границе доверия.

\vspace{0.25em}
\centering
\begin{tikzpicture}[
  font=\scriptsize,
  node distance=0.35cm,
  box/.style={rectangle, draw=KaspGreen!80!black, thick, rounded corners=2pt, align=center, inner sep=3pt, minimum height=0.85cm},
  arr/.style={-{Stealth[length=1.8mm]}, thick, KaspGreen}
]
\node[box, fill=KaspLight] (t) {TARA:\\угрозы, риски};
\node[box, right=0.45cm of t, fill=white] (g) {ЦБ\\(ГОСТ)};
\node[box, right=0.45cm of g, fill=white] (d) {Граница\\ДВБ};
\node[box, right=0.45cm of d, fill=KaspLight] (s) {SBOM:\\состав ПО};
\draw[arr] (t) -- (g);
\draw[arr] (g) -- (d);
\draw[arr] (d) -- node[above=1pt, font=\tiny] {обоснование} (s);
\end{tikzpicture}

\vspace{0.35em}
\raggedright
\textbf{Смысл:} без SBOM сложно проверить, что \textbf{реализованная} ДВБ совпадает с \textbf{спроектированной}; без TARA и целей SBOM остаётся «паспортом без адреса».
\end{frame}

\begin{frame}{ДВБ и граница доверия: зачем SBOM}
\small
\begin{itemize}
  \setlength{\itemsep}{0.2em}
  \item \textbf{ДВБ (TCB)} — подмножество системы, от которого зависит соблюдение политики безопасности; всё вне границы считается ненадёжным или внешним.
  \item SBOM позволяет: (1) перечислить компоненты внутри доверия; (2) отследить \textbf{транзитивные зависимости}, которые тоже входят в зону ответственности; (3) при угрозе на цепочку поставок — сопоставить CVE с \textbf{конкретной версией} в поставке.
  \item \textbf{Моделирование угроз} указывает, какие компоненты критичны для пути атаки; SBOM подтверждает их наличие и версии в артефакте.
\end{itemize}

\vspace{0.35em}
\centering
\begin{tikzpicture}[
  font=\scriptsize,
  comp/.style={rectangle, draw, rounded corners=2pt, inner sep=2pt, minimum height=0.55cm},
  arr/.style={-{Stealth[length=1.5mm]}, thick}
]
\node[ellipse, draw=KaspGreen, thick, minimum width=6.8cm, minimum height=2.8cm, fill=KaspLight] (zone) {};
\node[anchor=north west, font=\tiny] at (zone.north west) {Доверенная зона (концепция)};
\node[comp, fill=white] at (zone.center) {Компоненты по SBOM (библиотеки, модули, ОС-пакеты)};
\node[comp, fill=orange!15] at ([yshift=-0.55cm]zone.center) {Критичные элементы по TARA};
\end{tikzpicture}
\end{frame}

\begin{frame}{Совмещение с жизненным циклом (схема)}
\centering
\begin{tikzpicture}[
  font=\tiny,
  node distance=0.25cm,
  box/.style={rectangle, draw, align=center, inner sep=2pt, minimum height=0.62cm},
  arr/.style={-{Stealth[length=1.4mm]}, thick}
]
\node[box, fill=blue!8] (req) {Требования и\\ЦБ};
\node[box, right=0.4cm of req, fill=purple!8] (tara) {TARA};
\node[box, right=0.4cm of tara, fill=green!10] (arch) {Архитектура,\\ДВБ};
\node[box, right=0.4cm of arch, fill=yellow!10] (sbom) {SBOM в CI/CD};
\node[box, right=0.4cm of sbom, fill=red!8] (vex) {Уязвимости,\\VEX};
\draw[arr] (req) -- (tara);
\draw[arr] (tara) -- (arch);
\draw[arr] (arch) -- (sbom);
\draw[arr] (sbom) -- (vex);
\end{tikzpicture}

\vspace{0.45em}
\raggedright
\footnotesize
По §\,5.7--5.8 ГОСТ~Р~72118-2025 выбор конструктивных подходов должен \textbf{упрощать обоснование ЦБ}, в т.\,ч.\ за счёт учёта угроз, уязвимостей и снижения поверхности атаки; SBOM обеспечивает \textbf{проверяемый состав} ПО на выходе реализации и при анализе уязвимостей (VEX и др.).
\end{frame}

\section{Целеполагание SBOM и системная инженерия}

\begin{frame}{Анализ контекста и границы системы}
\footnotesize
\begin{itemize}
  \setlength{\itemsep}{0.18em}
  \item В \textbf{системной инженерии} (ЖЦ по ГОСТ~Р~57193 / ISO~15288) система определяется через \textbf{анализ контекста}: заинтересованные стороны, внешняя среда, границы, допущения и ограничения — это согласуется с \textbf{предположениями безопасности (ПБ)} в ГОСТ~Р~72118-2025.
  \item \textbf{Целеполагание SBOM:} границы учёта компонентов (продукт, сервис, образ, прошивка) задаются не «по умолчанию», а из контекста и архитектурного разбиения.
  \item Анализ контекста показывает, какие \textbf{домены} (данные, каналы, ОС, OTA, облако партнёра) нужно выделять для отдельных SBOM и отдельных ЦБ.
\end{itemize}
\end{frame}

\begin{frame}{ЦБ верхнего уровня и ЦБ подсистем / доменов}
\footnotesize
\begin{itemize}
  \setlength{\itemsep}{0.15em}
  \item \textbf{ЦБ верхнего уровня} формулируются для системы\/изделия в целом (например целостность критических функций, защита ПДн).
  \item \textbf{ЦБ подсистем и доменов безопасности} уточняют те же намерения для границ: бортовой шлюз, телематика, контейнер приложения, сегмент сети — с отдельным TARA\/риском по контуру.
  \item SBOM привязывают к \textbf{объекту учёта}: системный SBOM, SBOM сервиса, SBOM образа домена — с трассировкой к соответствующей ЦБ и к релизу.
\end{itemize}

\vspace{0.35em}
\centering
\begin{tikzpicture}[
  font=\tiny,
  box/.style={rectangle, draw=KaspGreen!75!black, rounded corners=2pt, align=center, inner sep=2pt, minimum height=0.55cm, text width=2.6cm}
]
\node[box, fill=KaspLight] (sys) {ЦБ системы\\(верхний уровень)};
\node[box, below left=0.45cm and -0.2cm of sys, fill=white] (a) {Домен A\\\scriptsize ЦБ$_A$ + SBOM$_A$};
\node[box, below=0.45cm of sys, fill=white] (b) {Домен B\\\scriptsize ЦБ$_B$ + SBOM$_B$};
\node[box, below right=0.45cm and -0.2cm of sys, fill=white] (c) {Домен C\\\scriptsize ЦБ$_C$ + SBOM$_C$};
\draw[->, thick, KaspGreen] (sys) -- (a);
\draw[->, thick, KaspGreen] (sys) -- (b);
\draw[->, thick, KaspGreen] (sys) -- (c);
\end{tikzpicture}
\end{frame}

\begin{frame}{Трассировка: контекст $\rightarrow$ ЦБ $\rightarrow$ политика SBOM}
\scriptsize
\setlength{\tabcolsep}{2.5pt}
\renewcommand{\arraystretch}{1.12}
\begin{tabularx}{\textwidth}{>{\bfseries}p{2.2cm}Y}
\toprule
Этап & Что фиксируют \\
\midrule
Контекст &
\makecell[tl]{Среда, интерфейсы, допущения;\\что в поставке и что внешнее} \\
ЦБ верхнего уровня &
\makecell[tl]{Намерения по ИБ для системы;\\вход к разбиению на домены} \\
ЦБ доменов &
\makecell[tl]{Уточнённые ЦБ и TARA по контуру;\\критичные компоненты} \\
Политика SBOM &
\makecell[tl]{Формат, глубина зависимостей, SBOM по доменам;\\привязка к сборке} \\
\bottomrule
\end{tabularx}

\vspace{0.35em}
\raggedright
\footnotesize
\textbf{Итог:} SBOM подтверждает \textbf{состав}, согласованный с инженерной декомпозицией и иерархией ЦБ; без этого состав «плавает» относительно модели угроз и ответственности.
\end{frame}

\section{Состав и идентификация компонентов}

\begin{frame}{Минимальный состав сведений о продукте}
\begin{table}
\footnotesize
\begin{tabularx}{\textwidth}{>{\bfseries}p{3.3cm}Y}
\toprule
Поле & Что должно быть указано \\
\midrule
Наименование продукта & Человеко-читаемое имя продукта, сервиса, пакета или образа. \\
Версия & Версия релиза, тега контейнера или поставки. \\
Поставщик / производитель & Организация-владелец или вендор. \\
Уникальный идентификатор & Внутренний ID, bom-ref, SPDXID или иной устойчивый идентификатор. \\
Дата формирования & Момент генерации SBOM. \\
Формат и версия схемы & Например, CycloneDX 1.5 или SPDX 2.3. \\
Связь с релизом & Коммит, build ID, digest контейнера, release tag. \\
\bottomrule
\end{tabularx}
\end{table}
\end{frame}

\begin{frame}{Минимальный состав сведений о компоненте}
\begin{table}
\scriptsize
\begin{tabularx}{\textwidth}{>{\bfseries}p{3.1cm}Y}
\toprule
Поле & Назначение \\
\midrule
Имя компонента & Название пакета или модуля. \\
Версия & Точная версия, а не «2.x». \\
Тип компонента & Библиотека, контейнер, ОС-пакет, модуль, приложение, фреймворк и т.д. \\
Поставщик / автор & Если известен. \\
Идентификатор & purl, CPE, SWID, внутренний идентификатор. \\
Контрольная сумма & Хэш бинарного артефакта или пакета. \\
Лицензия & SPDX-идентификатор, выражение лицензии, собственная лицензия. \\
Зависимости & Связи с родительскими и дочерними компонентами. \\
Статус включения & Прямая или транзитивная зависимость. \\
\bottomrule
\end{tabularx}
\end{table}
\end{frame}

\begin{frame}{Идентификация компонентов: что использовать и зачем}
\begin{table}
\tiny
\setlength{\tabcolsep}{2.5pt}
\renewcommand{\arraystretch}{1.08}
\begin{tabularx}{\textwidth}{>{\bfseries}p{1.65cm}p{3.35cm}Y}
\toprule
ID & Пример & Зачем \\
\midrule
purl &
\makecell[tl]{\texttt{pkg:maven/org.apache.}\\\texttt{logging.log4j/log4j-core@2.20.0}} &
Унифицированная идентификация пакетов для автоматизации. \\
CPE &
\makecell[tl]{\texttt{cpe:2.3:a:apache:log4j:}\\\texttt{2.20.0:*:*:*:*:*:*:*}} &
Сопоставление с частью баз уязвимостей. \\
SWID & Тег SWID & Учёт установленного ПО (ниши). \\
Хэш & SHA-256 файла или образа & Фиксация конкретного бинарного объекта. \\
Внутр.\,ID & \texttt{com.company.platform.auth-lib} & Собственные модули и форки. \\
\bottomrule
\end{tabularx}
\end{table}

\begin{block}{Практика}
Базовый набор для большинства организаций: purl + версия + хэш + внутренний идентификатор при необходимости.
\end{block}
\end{frame}

\begin{frame}{Отдельно учитывать внутренние и модифицированные компоненты}
\begin{itemize}
  \item собственные модули разработки;
  \item форки сторонних проектов;
  \item локально патченные библиотеки;
  \item переупакованные зависимости;
  \item компоненты, полученные от подрядчиков;
  \item компоненты без публичного реестра пакетов.
\end{itemize}

\begin{alertblock}{Типичная ошибка}
SBOM хорошо описывает open source-зависимости, но не отражает наиболее критичную внутреннюю часть продукта.
\end{alertblock}
\end{frame}

\begin{frame}{Политика включения зависимостей}
\begin{table}
\footnotesize
\begin{tabularx}{\textwidth}{>{\bfseries}p{3cm}Y}
\toprule
Категория & Рекомендация \\
\midrule
Прямые зависимости & Обязательны всегда. \\
Транзитивные зависимости & Обязательны, если реально входят в поставку или сборку. \\
Компоненты контейнерной базы & Обязательны для контейнерных поставок. \\
Системные пакеты ОС & Включать, если они поставляются вместе с образом или устройством. \\
Инструменты сборки & Обычно не включать в продуктовый SBOM, но можно учитывать в отдельном build-контуре. \\
Плагины, скрипты, модели & Включать, если они являются частью поставки или существенно влияют на поведение продукта. \\
\bottomrule
\end{tabularx}
\end{table}
\end{frame}

\section{Автоматизация и CI/CD}

\begin{frame}[fragile]{Диаграмма CI/CD: исходник PlantUML}
\textbf{Файл:} \texttt{images/sbom-cicd.puml}

\vspace{0.4em}

\lstinputlisting{images/sbom-cicd.puml}
\end{frame}

\begin{frame}{Диаграмма CI/CD: рендер и ссылка на PNG}
\begin{center}
\includegraphics[width=0.92\textwidth,height=0.72\textheight,keepaspectratio]{images/sbom-cicd.png}
\end{center}

\small
\textbf{PNG:} \href{run:images/sbom-cicd.png}{\texttt{images/sbom-cicd.png}}
\end{frame}

\begin{frame}{Как правильно встроить SBOM в CI/CD}
\begin{enumerate}
  \item Разработчик фиксирует зависимости и версии.
  \item Pipeline получает исходный код и сборочные метаданные.
  \item Генератор SBOM собирает данные из манифестов, lock-файлов, бинарных артефактов и контейнеров.
  \item Выполняется валидация схемы и контроль качества.
  \item SBOM подписывается и публикуется вместе с артефактом.
  \item Запускаются проверки безопасности, лицензий и политики поставки.
\end{enumerate}

\begin{block}{Лучший момент формирования}
Во время сборки релизного артефакта или контейнерного образа, а не постфактум вручную.
\end{block}
\end{frame}

\section{Качество, валидация и актуализация}

\begin{frame}{Контроль качества SBOM (1/2)}
\begin{table}
\scriptsize
\setlength{\tabcolsep}{2.5pt}
\begin{tabularx}{\textwidth}{>{\bfseries}p{2.9cm}Y}
\toprule
Проверка & Смысл \\
\midrule
Схема & CycloneDX/SPDX; проходит валидацию \\
Полнота & Нет явно пропущенных групп компонентов \\
Версии & У критичных — точные версии \\
Идентификаторы & purl, CPE, хэши, внутренние ID \\
\bottomrule
\end{tabularx}
\end{table}
\end{frame}

\begin{frame}{Контроль качества SBOM (2/2)}
\begin{table}
\scriptsize
\setlength{\tabcolsep}{2.5pt}
\begin{tabularx}{\textwidth}{>{\bfseries}p{2.9cm}Y}
\toprule
Проверка & Смысл \\
\midrule
Зависимости & Иерархия и связи корректны \\
Лицензии & Заполнены и нормализованы по возможности \\
Без дублей & Один компонент — однозначное описание \\
Прослеживаемость & Связь с build ID, образом, релизом \\
\bottomrule
\end{tabularx}
\end{table}
\end{frame}

\begin{frame}{Метрики качества SBOM (1/2)}
\begin{table}
\scriptsize
\begin{tabularx}{\textwidth}{>{\bfseries}p{3.1cm}Y}
\toprule
Метрика & Интерпретация \\
\midrule
Доля с точной версией & Пригодность для анализа уязвимостей \\
Доля с purl & Автосопоставление пакетов \\
Доля с лицензией & Лицензионный контроль \\
\bottomrule
\end{tabularx}
\end{table}
\end{frame}

\begin{frame}{Метрики качества SBOM (2/2)}
\begin{table}
\scriptsize
\begin{tabularx}{\textwidth}{>{\bfseries}p{3.1cm}Y}
\toprule
Метрика & Интерпретация \\
\midrule
Доля с хэшем & Подтверждение бинарных артефактов \\
Неидентифицированные & Ручной разбор и доработка процесса \\
Покрытие слоёв / ОС & Корректность контейнерных поставок \\
\bottomrule
\end{tabularx}
\end{table}
\end{frame}

\begin{frame}{Правила актуализации и версионности}
\begin{itemize}
  \item SBOM всегда должен быть привязан к конкретному состоянию продукта.
  \item Изменился состав поставки — требуется перевыпуск SBOM.
  \item Нужно хранить связки:
  \begin{itemize}
    \item SBOM $\leftrightarrow$ релиз;
    \item SBOM $\leftrightarrow$ коммит;
    \item SBOM $\leftrightarrow$ контейнерный digest;
    \item SBOM $\leftrightarrow$ поставка заказчику.
  \end{itemize}
  \item Предыдущие версии должны сохраняться для аудита и расследований.
\end{itemize}
\end{frame}

\section{Доверие, подписание и защита}

\begin{frame}[fragile]{Диаграмма доверия и публикации: исходник PlantUML}
\textbf{Файл:} \texttt{images/sbom-trust.puml}

\vspace{0.4em}

\lstinputlisting{images/sbom-trust.puml}
\end{frame}

\begin{frame}{Диаграмма доверия и публикации: рендер и ссылка на PNG}
\begin{center}
\includegraphics[width=0.92\textwidth,height=0.72\textheight,keepaspectratio]{images/sbom-trust.png}
\end{center}

\small
\textbf{PNG:} \href{run:images/sbom-trust.png}{\texttt{images/sbom-trust.png}}
\end{frame}
\begin{frame}{Почему SBOM нужно подписывать и защищать}
\begin{itemize}
  \item SBOM участвует в цепочке доверия поставки.
  \item Он должен быть защищен от подмены, потери или несанкционированного редактирования.
  \item Подпись и контроль целостности повышают доверие заказчика и аудитора.
  \item Хранилище SBOM должно обеспечивать:
  \begin{itemize}
    \item разграничение доступа;
    \item неизменяемость опубликованных релизных артефактов;
    \item журналирование операций;
    \item связь с релизными метаданными.
  \end{itemize}
\end{itemize}

\begin{block}{Минимум}
Подпись, хэш, доверенный репозиторий и трассировка, кто и когда сформировал документ.
\end{block}
\end{frame}

\section{Связь с уязвимостями и VEX}

\begin{frame}[fragile]{Диаграмма SBOM и VEX: исходник PlantUML}
\textbf{Файл:} \texttt{images/sbom-vex.puml}

\vspace{0.4em}

\lstinputlisting{images/sbom-vex.puml}
\end{frame}

\begin{frame}{Диаграмма SBOM и VEX: рендер и ссылка на PNG}
\begin{center}
\includegraphics[width=0.92\textwidth,height=0.72\textheight,keepaspectratio]{images/sbom-vex.png}
\end{center}

\small
\textbf{PNG:} \href{run:images/sbom-vex.png}{\texttt{images/sbom-vex.png}}
\end{frame}

\begin{frame}{Почему нельзя смешивать SBOM и выводы об эксплуатируемости}
\begin{columns}[T]
\begin{column}{0.48\textwidth}
\begin{block}{SBOM отвечает}
\begin{itemize}
  \item что входит в продукт;
  \item какие есть версии;
  \item каковы зависимости;
  \item какие идентификаторы и лицензии задействованы.
\end{itemize}
\end{block}
\end{column}
\begin{column}{0.48\textwidth}
\begin{block}{VEX / анализ уязвимостей отвечает}
\begin{itemize}
  \item затрагивает ли CVE конкретный продукт;
  \item достижима ли уязвимая функция;
  \item сработают ли компенсирующие меры;
  \item требуется ли обновление.
\end{itemize}
\end{block}
\end{column}
\end{columns}

\begin{alertblock}{Методический принцип}
SBOM — это состав. VEX — это интерпретация применимости уязвимости к продукту.
\end{alertblock}
\end{frame}

\section{Роли и методика}

\begin{frame}{Роли и ответственность}
\begin{table}
\scriptsize
\begin{tabularx}{\textwidth}{>{\bfseries}p{3cm}Y}
\toprule
Роль & Ответственность \\
\midrule
Разработка & Корректность описания зависимостей, внутренних модулей, форков, патчей. \\
DevOps / CI/CD & Автоматическая генерация, публикация, контроль прохождения pipeline. \\
Информационная безопасность & Требования к полноте, идентификации, качеству и использованию SBOM в анализе. \\
Релиз-менеджмент & Привязка SBOM к поставке, версии, образу и артефакту. \\
Юридическая функция / OSS Office & Лицензионная часть и обязательства по open source. \\
Архитектура / техлиды & Определение границ состава и правил включения компонентов. \\
\bottomrule
\end{tabularx}
\end{table}
\end{frame}

\begin{frame}{Как оформить внутреннюю методику по SBOM}
\begin{enumerate}
  \item Цель и область применения.
  \item Перечень продуктов и артефактов, где SBOM обязателен.
  \item Допустимые форматы: CycloneDX, SPDX.
  \item Обязательные поля и источники данных.
  \item Политика идентификации компонентов.
  \item Правила учета внутренних, внешних и модифицированных компонентов.
  \item Порядок автоматической генерации и публикации.
  \item Контроль качества, подпись, хранение и актуализация.
  \item Роли и ответственность.
\end{enumerate}
\end{frame}

\begin{frame}{Рекомендуемый «золотой минимум»}
\begin{itemize}
  \item Генерировать SBOM автоматически в CI/CD.
  \item Привязывать SBOM к каждому релизу и контейнерному образу.
  \item Включать прямые и транзитивные зависимости.
  \item Для контейнеров включать базовый образ и системные пакеты.
  \item Для каждого компонента фиксировать имя, версию, тип, purl, хэш, лицензию и связи зависимостей.
  \item Проводить валидацию схемы и полноты.
  \item Подписывать и хранить вместе с артефактом.
  \item Перевыпускать при изменении состава.
\end{itemize}
\end{frame}

\begin{frame}{Типичные ошибки внедрения}
\begin{itemize}
  \item SBOM формируется только «для галочки».
  \item Нет связи с конкретной сборкой или релизом.
  \item Учитываются только прямые зависимости.
  \item Игнорируются внутренние модули, форки и локальные патчи.
  \item Отсутствуют версии, purl, хэши и лицензии.
  \item Смешиваются состав продукта и итоговые выводы по уязвимостям.
  \item Процесс полностью ручной и не воспроизводимый.
  \item Нет политики по контейнерам, слоям и ОС-пакетам.
\end{itemize}
\end{frame}

\begin{frame}{Ландшафт инструментов (1/2)}
\begin{table}
\tiny
\setlength{\tabcolsep}{2.5pt}
\begin{tabularx}{\textwidth}{>{\bfseries}p{2.45cm}Y Y}
\toprule
Класс & Примеры & Роль \\
\midrule
Генерация SBOM &
Syft, cdxgen, CycloneDX plugins, SPDX &
\makecell[tl]{Документ из исходников,\\пакетов, контейнеров} \\
Сканирование &
\makecell[tl]{OWASP Dependency-Track,\\Trivy, Grype, Snyk} &
\makecell[tl]{Состав; известные\\уязвимости} \\
Лицензии &
FOSSA, ScanCode, OSS-review &
\makecell[tl]{Лицензии и\\обязательства} \\
\bottomrule
\end{tabularx}
\end{table}
\end{frame}

\begin{frame}{Ландшафт инструментов (2/2)}
\begin{table}
\footnotesize
\begin{tabularx}{\textwidth}{>{\bfseries}p{2.6cm}Y Y}
\toprule
Класс & Примеры & Роль \\
\midrule
CI/CD &
GitLab CI, Jenkins, GitHub Actions, registry &
\makecell[tl]{Генерация, публикация,\\хранение} \\
Подпись &
cosign, in-toto, внутренний PKI &
Доверие и целостность \\
\bottomrule
\end{tabularx}
\end{table}
\end{frame}

\section{Дорожная карта внедрения}

\begin{frame}[fragile]{Диаграмма поэтапного внедрения: исходник PlantUML}
\textbf{Файл:} \texttt{images/sbom-rollout.puml}

\vspace{0.4em}

\lstinputlisting{images/sbom-rollout.puml}
\end{frame}

\begin{frame}{Диаграмма поэтапного внедрения: рендер и ссылка на PNG}
\begin{center}
\includegraphics[width=0.92\textwidth,height=0.72\textheight,keepaspectratio]{images/sbom-rollout.png}
\end{center}

\small
\textbf{PNG:} \href{run:images/sbom-rollout.png}{\texttt{images/sbom-rollout.png}}
\end{frame}

\begin{frame}{Практический план внедрения: этапы 1--2}
\begin{table}
\scriptsize
\setlength{\tabcolsep}{2.5pt}
\renewcommand{\arraystretch}{1.1}
\begin{tabularx}{\textwidth}{>{\bfseries}p{1.35cm}Y Y}
\toprule
№ & Что делаем & Результат \\
\midrule
1 &
\makecell[tl]{Цели, форматы, продукты,\\пилотная команда} &
Утверждённый минимальный стандарт SBOM \\
2 &
\makecell[tl]{Генерация в CI/CD и\\контейнерном pipeline} &
Повторяемый выпуск SBOM для релизов \\
\bottomrule
\end{tabularx}
\end{table}
\end{frame}

\begin{frame}{Практический план внедрения: этапы 3--4}
\begin{table}
\scriptsize
\setlength{\tabcolsep}{2.5pt}
\renewcommand{\arraystretch}{1.1}
\begin{tabularx}{\textwidth}{>{\bfseries}p{1.35cm}Y Y}
\toprule
№ & Что делаем & Результат \\
\midrule
3 &
\makecell[tl]{Валидация, метрики, подпись,\\хранение} &
Управляемый и проверяемый процесс \\
4 &
\makecell[tl]{Охват, VEX, аудит,\\заказчики} &
SBOM — штатный артефакт поставки \\
\bottomrule
\end{tabularx}
\end{table}
\end{frame}

\section{Итоги}

\begin{frame}{Ключевые выводы (1/2)}
\small
\begin{itemize}
  \setlength{\itemsep}{0.2em}
  \item Методически правильный SBOM — регламентированный артефакт ЖЦ ПО.
  \item С \textbf{ГОСТ~Р~72118-2025} SBOM даёт \textbf{доказуемый состав} для аудита и рисков.
  \item \textbf{TARA} и \textbf{ЦБ} задают, что защищаем; \textbf{ДВБ} — границу доверия; SBOM — фактический состав в контуре поставки.
  \item Целеполагание SBOM согласуется с системной инженерией: контекст $\rightarrow$ ЦБ системы и доменов $\rightarrow$ политика SBOM.
\end{itemize}
\end{frame}

\begin{frame}{Ключевые выводы (2/2)}
\small
\begin{itemize}
  \setlength{\itemsep}{0.2em}
  \item Качество: границы состава, идентификаторы, автоматизация.
  \item Нужны несколько источников: манифесты, сборка, бинарный и контейнерный анализ.
  \item SBOM — актуальный, проверяемый, подписанный, привязанный к релизу.
  \item Применимость уязвимостей — отдельный слой (например VEX), не смешивать с составом.
\end{itemize}
\end{frame}

\begin{frame}{Короткая формула}
\begin{block}{Определение}
Методически правильный SBOM — это формализованный, автоматизированный, проверяемый и актуальный перечень компонентов конкретного программного продукта или поставки, встроенный в жизненный цикл разработки и поставки.
\end{block}

\vspace{1em}

\begin{center}
\Large Спасибо за внимание
\end{center}
\end{frame}

\begin{frame}
\begin{center}
\Huge Вопросы?
\end{center}
\end{frame}

\end{document}
